\section{Proofs}
%%%
\subsection{Proof of \cref{lem:upper_weight}}~\label{sec:upper_weight}
\begin{proof}
    By \cref{lem:taylor} and the assumption that $\mu\left( \bigcup_{\tau = 0}^{t}\tx(\tau)\Delta \bigcup_{\tau = 0}^{t+\Delta t}\tx(\tau)\right)\geq \delta,$ we have
    \begin{align*}
       E\left(\bigcup_{\tau = 0}^{t}\tx(\tau)\right) & = E\left(\bigcup_{\tau = 0}^{t+\Delta t}\tx(\tau)\right) - \Delta t\Big[\mu\left( \bigcup_{\tau = 0}^{t}\tx(\tau)\Delta \bigcup_{\tau = 0}^{t+\Delta t}\tx(\tau)\right)\cdot \displaystyle\int_{\bigcup_{\tau = 0}^{t}\tx(\tau)} \ell(\hat{f}(\weight,\psi))d\mu\\
        & +\displaystyle\int_{\bigcup_{\tau = 0}^{t}\tx(\tau)} \ell'(\hat{f}(\weight,\psi))\cdot \partial_{\weight}^1 \hat{f}(\weight,\psi)\cdot \Delta w \hspace{1mm}  d\mu\Big]- o(\Delta t)\\
        & = E\left(\bigcup_{\tau = 0}^{t+\Delta t}\tx(\tau)\right) + \Delta t\left(-\mu\left( \bigcup_{\tau = 0}^{t}\tx(\tau)\Delta \bigcup_{\tau = 0}^{t+\Delta t}\tx(\tau)\right)\right)\cdot \displaystyle\int_{\bigcup_{\tau = 0}^{t}\tx(\tau)\Delta \bigcup_{\tau = 0}^{t+\Delta t}\tx(\tau)} \ell(\hat{f}(\weight,\psi))d\mu\\
        & - \Delta t\displaystyle\int_{\bigcup_{\tau = 0}^{t}\tx(\tau)} \ell'(\hat{f}(\weight,\psi))\cdot \partial_{\weight}^1 \hat{f}(\weight,\psi)\cdot \Delta w \hspace{1mm}  d\mu- o(\Delta t)\\
        & \leq E\left(\bigcup_{\tau = 0}^{t+\Delta t}\tx(\tau)\right) + \Delta t \cdot \delta \cdot \displaystyle\int_{\bigcup_{\tau = 0}^{t}\tx(\tau)\Delta \bigcup_{\tau = 0}^{t+\Delta t}\tx(\tau)} \ell(\hat{f}(\weight,\psi))d\mu\\
        & - \Delta t\displaystyle\int_{\bigcup_{\tau = 0}^{t}\tx(\tau)} \ell'(\hat{f}(\weight,\psi))\cdot \partial_{\weight}^1 \hat{f}(\weight,\psi)\cdot \Delta w \hspace{1mm}  d\mu- o(\Delta t)\\
    \end{align*}
    Subtracting $E\left(\bigcup_{\tau = 0}^{t+\Delta t}\tx(\tau)\right)$ from both sides and combining like terms,
    \begin{align*}
         E\left(\bigcup_{\tau = 0}^{t}\tx(\tau)\right) - E\left(\bigcup_{\tau = 0}^{t+\Delta t}\tx(\tau)\right)
         & \leq \Delta t \cdot \delta \cdot \displaystyle\int_{\bigcup_{\tau = 0}^{t}\tx(\tau)\Delta \bigcup_{\tau = 0}^{t+\Delta t}\tx(\tau)} \ell(\hat{f}(\weight,\psi))d\mu\\
        & - \Delta t\displaystyle\int_{\bigcup_{\tau = 0}^{t}\tx(\tau)} \ell'(\hat{f}(\weight,\psi))\cdot \partial_{\weight}^1 \hat{f}(\weight,\psi)\cdot \Delta w \hspace{1mm}  d\mu- o(\Delta t).
    \end{align*}
    Additionally, we can assume that the loss function $\ell$ is bounded above by a constant $M_0,$ so
    \begin{align*}
        E\left(\bigcup_{\tau = 0}^{t}\tx(\tau)\right) - E\left(\bigcup_{\tau = 0}^{t+\Delta t}\tx(\tau)\right)
        & \leq \Delta t \cdot \delta \cdot M_0- \Delta t\displaystyle\int_{\bigcup_{\tau = 0}^{t}\tx(\tau)} \ell'(\hat{f}(\weight,\psi))\cdot \partial_{\weight}^1 \hat{f}(\weight,\psi)\cdot \Delta w \hspace{1mm}  d\mu- o(\Delta t).
    \end{align*}
    To consider this difference in expected values across future tasks, set $\Delta t =1 $ and let us sum across such future tasks, 
    \begin{align*}
         \sum_{\tau = t}^{T} \left(E\left(\bigcup_{\nu = 0}^{\tau}\tx(\nu)\right) - E\left(\bigcup_{\nu = 0}^{\tau+1}\tx(\nu)\right)\right)
        & \leq \sum_{\tau =t}^{T} \Big( M_0\cdot \delta - \displaystyle\int_{\bigcup_{\nu = 0}^{\tau}\tx(\nu)} \ell'(\hat{f}(\weight,\psi))\cdot \partial_{\weight}^1 \hat{f}(\weight,\psi)\cdot \Delta w \hspace{1mm}  d\mu \Big).
    \end{align*}
    Now, we can choose the following for each $\tau$
    \begin{align*}
        \Delta \weight & =  \ell'(\hat{f}(\weight,\psi))\cdot \partial_{\weight}^1 \hat{f}(\weight,\psi).
    \end{align*}
    Thus,
    \begin{align*}
         \sum_{\tau = t}^{T} \left(E\left(\bigcup_{\nu = 0}^{\tau}\tx(\nu)\right) - E\left(\bigcup_{\nu = 0}^{\tau+1}\tx(\nu)\right)\right)
        & \leq \sum_{\tau =t}^{T} \Big( M_0\cdot \delta - \displaystyle\int_{\bigcup_{\nu = 0}^\tau\tx(\nu)} \left(\ell'(\hat{f}(\weight,\psi))\cdot \partial_{\weight}^1 \hat{f}(\weight,\psi)\right)^2 \hspace{1mm}  d\mu \Big).
    \end{align*}
    Now, notice the following
    \begin{align*}
        J(\weight(\tau),\psi(\tau),\tx(\tau)) & = \int_{\bigcup_{\nu = 0}^\tau \tx(\nu)} \ell(\hat{f}(\weight,\psi))d\mu = E\left(\bigcup_{\nu = 0}^{\tau}\tx(\nu)\right).
    \end{align*}
    Thus,
    \begin{align*}
         \sum_{\tau = t}^{T} \left(J(\tx(\tau)) - J(\tx(\tau+1))\right)
         & = \sum_{\tau = t}^{T} \left(E\left(\bigcup_{\nu = 0}^{\tau}\tx(\nu)\right) - E\left(\bigcup_{\nu = 0}^{\tau+1}\tx(\nu)\right)\right)\\
        & \leq \sum_{\tau =t}^{T} \Big( M_0\cdot \delta - \displaystyle\int_{\bigcup_{\nu = 0}^\tau\tx(\nu)} \left(\ell'(\hat{f}(\weight,\psi))\cdot \partial_{\weight}^1 \hat{f}(\weight,\psi)\right)^2 \hspace{1mm}  d\mu \Big),
    \end{align*}
    as desired.
\end{proof}



\subsection{Proof of \cref{lem:upper_arch}}~\label{sec:upper_arch}
\begin{proof}
    By \cref{lem:taylor} and the assumption that $\mu\left( \bigcup_{\tau = 0}^{t}\tx(\tau)\Delta \bigcup_{\tau = 0}^{t+\Delta t}\tx(\tau)\right)\geq \delta,$ we have
    \begin{align*}
       E\left(\bigcup_{\tau = 0}^{t}\tx(\tau)\right) & = E\left(\bigcup_{\tau = 0}^{t+\Delta t}\tx(\tau)\right) - \Delta t\Big[\mu\left( \bigcup_{\tau = 0}^{t}\tx(\tau)\Delta \bigcup_{\tau = 0}^{t+\Delta t}\tx(\tau)\right)\cdot \displaystyle\int_{\bigcup_{\tau = 0}^{t}\tx(\tau)} \ell(\hat{f}(\weight,\psi))d\mu\\
        & +\displaystyle\int_{\bigcup_{\tau = 0}^{t}\tx(\tau)} \ell'(\hat{f}(\weight,\psi))\cdot \partial_{\weight}^1 \hat{f}(\weight,\psi)\cdot \Delta w \hspace{1mm}  d\mu\\
        & +\displaystyle\int_{\bigcup_{\tau = 0}^{t}\tx(\tau)} \ell'(\hat{f}(\weight,\psi))\cdot \partial_{\psi}^1 \hat{f}(\weight,\psi)\cdot \Delta \psi \hspace{1mm}  d\mu\Big]- o(\Delta t)\\
        & = E\left(\bigcup_{\tau = 0}^{t+\Delta t}\tx(\tau)\right) + \Delta t\left(-\mu\left( \bigcup_{\tau = 0}^{t}\tx(\tau)\Delta \bigcup_{\tau = 0}^{t+\Delta t}\tx(\tau)\right)\right)\cdot \displaystyle\int_{\bigcup_{\tau = 0}^{t}\tx(\tau)\Delta \bigcup_{\tau = 0}^{t+\Delta t}\tx(\tau)} \ell(\hat{f}(\weight,\psi))d\mu\\
        & - \Delta t\displaystyle\int_{\bigcup_{\tau = 0}^{t}\tx(\tau)} \ell'(\hat{f}(\weight,\psi))\cdot \partial_{\weight}^1 \hat{f}(\weight,\psi)\cdot \Delta w \hspace{1mm}  d\mu\\
        &-\Delta t \displaystyle\int_{\bigcup_{\tau = 0}^{t}\tx(\tau)} \ell'(\hat{f}(\weight,\psi))\cdot \partial_{\psi}^1 \hat{f}(\weight,\psi)\cdot \Delta \psi \hspace{1mm}  d\mu- o(\Delta t)\\
        & \leq E\left(\bigcup_{\tau = 0}^{t+\Delta t}\tx(\tau)\right) + \Delta t \cdot \delta \cdot \displaystyle\int_{\bigcup_{\tau = 0}^{t}\tx(\tau)\Delta \bigcup_{\tau = 0}^{t+\Delta t}\tx(\tau)} \ell(\hat{f}(\weight,\psi))d\mu\\
        & - \Delta t\displaystyle\int_{\bigcup_{\tau = 0}^{t}\tx(\tau)} \ell'(\hat{f}(\weight,\psi))\cdot \partial_{\weight}^1 \hat{f}(\weight,\psi)\cdot \Delta w \hspace{1mm}  d\mu\\
        & -\Delta t \displaystyle\int_{\bigcup_{\tau = 0}^{t}\tx(\tau)} \ell'(\hat{f}(\weight,\psi))\cdot \partial_{\psi}^1 \hat{f}(\weight,\psi)\cdot \Delta \psi \hspace{1mm}  d\mu- o(\Delta t).
    \end{align*}
    Subtracting $E\left(\bigcup_{\tau = 0}^{t+\Delta t}\tx(\tau)\right)$ from both sides and combining like terms,
    \begin{align*}
         E\left(\bigcup_{\tau = 0}^{t}\tx(\tau)\right) - E\left(\bigcup_{\tau = 0}^{t+\Delta t}\tx(\tau)\right)
         & \leq \Delta t \cdot \delta \cdot \displaystyle\int_{\bigcup_{\tau = 0}^{t}\tx(\tau)\Delta \bigcup_{\tau = 0}^{t+\Delta t}\tx(\tau)} \ell(\hat{f}(\weight,\psi))d\mu\\
        & - \Delta t\displaystyle\int_{\bigcup_{\tau = 0}^{t}\tx(\tau)} \ell'(\hat{f}(\weight,\psi))\cdot \partial_{\weight}^1 \hat{f}(\weight,\psi)\cdot \Delta w \hspace{1mm}  d\mu\\
        & -\Delta t \displaystyle\int_{\bigcup_{\tau = 0}^{t}\tx(\tau)} \ell'(\hat{f}(\weight,\psi))\cdot \partial_{\psi}^1 \hat{f}(\weight,\psi)\cdot \Delta \psi \hspace{1mm}  d\mu- o(\Delta t).
    \end{align*}
    Additionally, we can assume that the loss function $\ell$ is bounded above by a constant $M_0,$ so
    \begin{align*}
        E\left(\bigcup_{\tau = 0}^{t}\tx(\tau)\right) - E\left(\bigcup_{\tau = 0}^{t+\Delta t}\tx(\tau)\right)
        & \leq \Delta t \cdot \delta \cdot M_0- \Delta t\displaystyle\int_{\bigcup_{\tau = 0}^{t}\tx(\tau)} \ell'(\hat{f}(\weight,\psi))\cdot \partial_{\weight}^1 \hat{f}(\weight,\psi)\cdot \Delta w \hspace{1mm}  d\mu\\
        & -\Delta t \displaystyle\int_{\bigcup_{\tau = 0}^{t}\tx(\tau)} \ell'(\hat{f}(\weight,\psi))\cdot \partial_{\psi}^1 \hat{f}(\weight,\psi)\cdot \Delta \psi \hspace{1mm}  d\mu- o(\Delta t).
    \end{align*}
    To consider this difference in expected values across future tasks, set $\Delta t =1 $ and let us sum across such future tasks, 
    \begin{align*}
         \sum_{\tau = t}^{T} \left(E\left(\bigcup_{\nu = 0}^{\tau}\tx(\nu)\right) - E\left(\bigcup_{\nu = 0}^{\tau+1}\tx(\nu)\right)\right)
        & \leq \sum_{\tau =t}^{T} \Big( M_0\cdot \delta - \displaystyle\int_{\bigcup_{\nu = 0}^{\tau}\tx(\nu)} \ell'(\hat{f}(\weight,\psi))\cdot \partial_{\weight}^1 \hat{f}(\weight,\psi)\cdot \Delta w \hspace{1mm}  d\mu\\
        & -\displaystyle\int_{\bigcup_{\tau = 0}^{t}\tx(\tau)} \ell'(\hat{f}(\weight,\psi))\cdot \partial_{\psi}^1 \hat{f}(\weight,\psi)\cdot \Delta \psi \hspace{1mm}  d\mu\Big).
    \end{align*}
    Now, we can choose the following for each $\nu$
    \begin{align*}
        \Delta \weight & =  \ell'(\hat{f}(\weight,\psi))\cdot \partial_{\weight}^1 \hat{f}(\weight,\psi).
    \end{align*}
    and
    \begin{align*}
        \Delta \psi & =  \ell'(\hat{f}(\weight,\psi))\cdot \partial_{\psi}^1 \hat{f}(\weight,\psi).
    \end{align*}
    Thus,
    \begin{align*}
         \sum_{\tau = t}^{T} \left(E\left(\bigcup_{\nu = 0}^{\tau}\tx(\nu)\right) - E\left(\bigcup_{\nu = 0}^{\tau+1}\tx(\nu)\right)\right)
        & \leq \sum_{\tau =t}^{T} \Big( M_0\cdot \delta - \displaystyle\int_{\bigcup_{\nu = 0}^\tau\tx(\nu)} \left(\ell'(\hat{f}(\weight,\psi))\cdot \partial_{\weight}^1 \hat{f}(\weight,\psi)\right)^2 \hspace{1mm}  d\mu\\
        & -\displaystyle\int_{\bigcup_{\nu = 0}^\tau\tx(\nu)} \left(\ell'(\hat{f}(\weight,\psi))\cdot \partial_{\psi}^1 \hat{f}(\weight,\psi)\right)^2 \hspace{1mm}  d\mu \Big).
    \end{align*}
    Now, notice the following
    \begin{align*}
        J(\weight(\tau),\psi(\tau),\tx(\tau)) & = \int_{\bigcup_{\nu = 0}^\tau \tx(\nu)} \ell(\hat{f}(\weight,\psi))d\mu = E\left(\bigcup_{\nu = 0}^{\tau}\tx(\nu)\right).
    \end{align*}
    Thus,
    \begin{align*}
         \sum_{\tau = t}^{T} \left(J(\tx(\tau)) - J(\tx(\tau+1))\right)
         & = \sum_{\tau = t}^{T} \left(E\left(\bigcup_{\nu = 0}^{\tau}\tx(\nu)\right) - E\left(\bigcup_{\nu = 0}^{\tau+1}\tx(\nu)\right)\right)\\
        & \leq \sum_{\tau =t}^{T} \Big( M_0\cdot \delta - \displaystyle\int_{\bigcup_{\nu = 0}^\tau\tx(\nu)} \left(\ell'(\hat{f}(\weight,\psi))\cdot \partial_{\weight}^1 \hat{f}(\weight,\psi)\right)^2 \hspace{1mm}  d\mu\\
        & -\displaystyle\int_{\bigcup_{\nu = 0}^\tau\tx(\nu)} \left(\ell'(\hat{f}(\weight,\psi))\cdot \partial_{\psi}^1 \hat{f}(\weight,\psi)\right)^2 \hspace{1mm}  d\mu \Big)\\
    \end{align*}
    as desired.
\end{proof}

\subsection{Proof of \cref{HJB}}\label{sec:HJB}
\begin{proof}
    Let $J, V,$ and $V^*$ be as defined above. To begin, we split the sum in $V(t)$ over the discrete intervals $[t,t+\Delta t]$ and $[t+\Delta t, T].$ Observe,
    \begin{align}\label{b}
        V^*(t) & = \min_{w\in \mathcal{W}(\psi^*(t))} \int_{t}^T J(w(\tau),\psi^*(t), \tx(\tau))d\tau\nonumber\\
                & = \min_{w\in \mathcal{W}(\psi^*(t))} \biggr[\int_{t}^{t+\Delta t} J(w(\tau),\psi^*(t), \tx(\tau))d\tau+ \int_{t+\Delta t}^{T} J(w(\tau),\psi^*(t), \tx(\tau))d\tau\biggr]\nonumber\\
                & = \min_{w\in \mathcal{W}(\psi^*(t))} \int_{t}^{t+\Delta t} J(w(\tau),\psi^*(t),\tx(\tau))d\tau + V^*(t+\Delta t)\nonumber\\
                & = \min_{w\in \mathcal{W}(\psi^*(t))} J(\weight(t),\psi^*(t),\tx(t))\Delta t + V^*(t+\Delta t).
    \end{align}
    Now, we provide the Taylor series expansion of $V^*(t+\Delta t)$ about $t.$ Notice,
    \begin{align}
        V^*(t+\Delta t) & = V^*(t) + \Delta t \left[\partials_{t}V^*(t)  + \partials_{\tx}V^*(t) d_{t} \tx + \partials_{\weight}V^*(t)  \partials_t \weight \right] + o(\Delta t)\nonumber\\
                & = V^*(t) + \Delta t\biggr[\partials_{t} V^*(t)  + \partials_{\tx}V^*(t) d_{t} \tx + \partials_{\weight}V^*(t) [\mathrm{A}^*(t)\weight(t)( \mathrm{B}^*(t))^T + u(t)]\biggr] + o(\Delta t),\label{a}
    \end{align}
    where $u(t)$ represents the updates made to the each weights matrix of the new dimensions.
    Substituting \eqref{a} into \eqref{b}, we have
    \begin{align*}
        V^*(t) & = \min_{w\in \mathcal{W}(\psi^*(t))} J(\weight(t),\psi^*(t),\tx(t))\Delta t+ V^*(t) + \Delta t\biggr[\partials_{t} V^*(t)  + \partials_{\tx}V^*(t) d_{t} \tx\\ 
                &+ \partials_{\weight}V^*(t) [\mathrm{A}^*(t)\weight(t)( \mathrm{B}^*(t))^T + u(t)]\biggr] + o(\Delta t).
    \end{align*}
    Canceling $V^*(t)$ gives
    \begin{align*}
        0 & = \min_{w\in \mathcal{W}(\psi^*(t))} J(\weight(t),\psi^*(t),\tx(t))\Delta t+ \Delta t\left[\partials_{t}V^*(t)  + \partials_{\tx}V^*(t) d_{t} \tx \right.  \\ &+
        \left. \partials_{\weight}V^*(t) [\mathrm{A}^*(t)\weight(t)( \mathrm{B}^*(t))^T + u(t)]\right] + o(\Delta t).
    \end{align*}
    Dividing both sides by $\Delta t$ produces
    \begin{align*}
        0 & = \min_{w\in \mathcal{W}(\psi^*(t))} J(\weight(t),\psi^*(t),\tx(t)) +
        \partials_{t}V^*(t)  + \partials_{\tx}V^*(t) V^*(t) d_{t} \tx + \partials_{\weight}V^*(t) [\mathrm{A}^*(t)\weight(t)( \mathrm{B}^*(t))^T + u(t)].
    \end{align*}
    Finally, reordering gives
    \begin{align*}
        -\partials_{t} V^*(t) & = \min_{w\in \mathcal{W}(\psi^*(t))} J(\weight(t),\psi^*(t),\tx(t))+ \partials_{\tx}V^*(t) d_{t} \tx + \partials_{\weight}V^*(t) V^*(t) [\mathrm{A}^*(t)\weight(t)( \mathrm{B}^*(t))^T + u(t)],
    \end{align*}
    as desired.
\end{proof}

\subsection{Proof of \cref{thm:lower_HJB}}\label{sec:lower_HJB}
\begin{proof}
Given \cref{HJB}, assume there exists an stochastic gradient based optimization procedure  and choose $ u(t) =- \sum^{I} \alpha(i) \mathrm{g}^{(i)}$  such that $\min_{\weight \in \mathcal{W}(\psi^*(t))} J(\weight(t),\psi^*(t),\tx(t)) \leq \varepsilon,$ where $I$ is the number of updates. Then, we have 
\begin{align*}
    -\partials_{t}V^*(t) & \leq \varepsilon + \partials_{\tx}V^*(t) d_{t} \tx +  \partials_{\weight}V^*(t)  \left( \mathrm{A}^*(t)\weight(t)( \mathrm{B}^*(t))^T - \sum^{I} \alpha(i) \mathrm{g}^{(i)}  \right)  \\    
\end{align*}
For any $dt$ in terms of tasks $t$  let $\partials_{\tx}V^*(t) d_{t} \tx \leq \norm{\partials_{\tx}V^*(t) } \norm{d_{t} \tx} \leq \sing \delta$ as $\norm{d_{t} \tx}  \geq \mu(\tx(t)\bigtriangleup \tx(t+1)) \geq \delta$ and $\norm{\partials_{\tx}V^*(t) }$ is less than the largest singular value of $\partials_{\tx}V^*(t) .$  Thus we write 
\begin{align*}
    -\partials_{t}V^*(t) & \leq \varepsilon + \sing \delta +  \partials_{\weight}V^*(t)  \left( \mathrm{A}^*(t)\weight(t)( \mathrm{B}^*(t))^T - \sum^{I} \alpha(i) \mathrm{g}^{(i)}  \right) 
\end{align*}
Taking $g_{MIN}$ to be the smallest value of the gradient over all update iterations, we have 
\begin{align*}
    -\partials_{t} & \leq \varepsilon + \sing \delta +  \partials_{\weight}V^*(t)  \left( \mathrm{A}^*(t)\weight(t) (\mathrm{B}^*(t))^T - \mathrm{g}_{\mathrm{MIN} } \sum^{I} \alpha(i)  \right) \\
     & \leq \varepsilon + \sing \delta +  \norm{\partials_{\weight}V^*(t) } \norm{\mathrm{A}^*(t)\weight(t)( \mathrm{B}^*(t))^T} - \partials_{\weight}V^*(t)  \left(  \mathrm{g}_{\mathrm{MIN} } \sum^{I} \alpha(i)  \right) \\
     & \leq \varepsilon + \left[\sing \delta +  \singw \norm{\mathrm{A}^*(t)\weight(t)( \mathrm{B}^*(t))^T} - \singmin \norm{\mathrm{g}_{\mathrm{MIN} }} \norm{\sum^{I} \alpha(i) } \right]
\end{align*}
This finally provides 
\begin{align*}
    \partials_{t} & \leq \mathrm{sup} \varepsilon - \mathrm{inf} \left[\sing \delta +  \singw \norm{\mathrm{A}^*(t)\weight(t)( \mathrm{B}^*(t))^T} - \singmin \norm{\mathrm{g}_{\mathrm{MIN} }} \norm{\sum^{I} \alpha(i) } \right]
\end{align*}
For the total variation to be upper bounded, we need the quantity in the brackets to go to zero, this provides.
\begin{align*}
\left[\sing \delta +  \singw \norm{\mathrm{A}^*(t)\weight(t)( \mathrm{B}^*(t))^T} - \singmin \norm{\mathrm{g}_{\mathrm{MIN} }} \norm{\sum^{I} \alpha(i) } \right] = 0 \\
\end{align*}
\begin{align*}
   \sing \delta &= \left[ \singmin \norm{\mathrm{g}_{\mathrm{MIN} }}  \norm{\sum^{I} \alpha(i) } - \singw \norm{\mathrm{A}^*(t)\weight(t)( \mathrm{B}^*(t))^T} \right]
\end{align*}
providing the result
\end{proof}

